\documentclass[11pt]{article}
\usepackage[utf8]{inputenc}
\usepackage[T1]{fontenc}
\usepackage[spanish,es-noshorthands]{babel}
\usepackage[margin=2.3cm]{geometry}
\usepackage{amsmath,amssymb}
\usepackage{enumitem}
\setlength{\parindent}{0pt}
\setlength{\parskip}{6pt}

\begin{document}

\begin{center}
{\Large\bfseries Control 1 --- Microeconomía II}\\[4pt]
{\large Versión «lo más probable» (según la predicción del 21-10)}\\[4pt]
80 puntos \quad$\cdot$\quad 80 minutos \quad$\cdot$\quad sin apuntes
\end{center}

\hrule
\medskip

\textbf{I. Comentes (5 pts c/u).} Comente si las siguientes afirmaciones son verdaderas, falsas o inciertas. Justifique su respuesta; la nota depende de la justificación.

\begin{enumerate}[leftmargin=*, itemsep=8pt]
\item Un monopolista maximizador de beneficios puede elegir operar en el tramo inelástico de la demanda si allí el precio es suficientemente alto.

\item La pérdida irrecuperable de eficiencia del monopolio corresponde a toda la pérdida de excedente del consumidor.

\item Cuando el monopolista puede aplicar discriminación perfecta, la pérdida de eficiencia respecto del mercado competitivo es mayor que con monopolio de precio único.

\item En el modelo de Hotelling con dos empresas y precios iguales, la competencia lleva a las empresas a ubicarse en la localización socialmente óptima.
\end{enumerate}

\textbf{II. Desarrollo (10 pts c/u).}

\begin{enumerate}[leftmargin=*, itemsep=8pt, resume]
\item Explique la relación entre la elasticidad de la demanda y el poder de mercado del monopolista. ¿Qué ocurre con ese poder cuando la demanda se vuelve perfectamente elástica? (Pista: Lerner.)

\item Explique por qué el monopolio no es eficiente en sentido de Pareto. Grafique.
\end{enumerate}

\textbf{III. Matemático (40 pts).}

Un monopolista posee una función de costos totales de la forma $c(q) = q^{2} + 20q + 300$ y abastece a un mercado cuya función de demanda es $q(p) = 100 - \dfrac{p}{3}$.

\begin{enumerate}[label=\alph*), leftmargin=*, itemsep=8pt]
\item \textbf{(20 pts)} Obtenga todos los elementos del equilibrio del monopolista.

\item \textbf{(20 pts)} Calcule la pérdida de eficiencia que sufre esta economía. En específico, señale cómo varían el precio y la cantidad de equilibrio y los excedentes del consumidor, del productor y social respecto de la situación de competencia perfecta. ¿Produce el monopolista la mitad de la cantidad de competencia perfecta? Explique. Grafique ambos apartados.
\end{enumerate}

\medskip
\hrule
\smallskip
{\small\textit{Antes de empezar, los 4 chequeos:} (1) reparte el tiempo por puntaje (matemático 40 min, desarrollos 10 c/u, comentes 5 c/u); (2) subraya el sujeto de cada pregunta; (3) cierra cada «compare» con la comparación explícita; (4) no reuses cifras de otro ítem sin revisarlas. Veredicto V/F/I en la primera palabra; trabaja con fracciones; $\pi$ con paréntesis; $\varepsilon$ evaluada en el punto.}

\end{document}
